% =============================================================================
\section{Project Overview}
% =============================================================================

\subsection{Purpose}

The \texttt{star-proto} repository contains shared Protocol Buffer (protobuf) definitions used across the STAR robot system:

\begin{itemize}
    \item \textbf{star-gateway}: Kotlin/Java gRPC server
    \item \textbf{star-ui}: TypeScript gRPC-Web client
    \item \textbf{star-esp32-firmware}: (Future) ESP32 binary protocol
\end{itemize}

\subsection{Why Protocol Buffers?}

\begin{table}[H]
\centering
\begin{tabular}{|l|l|l|}
\hline
\textbf{Feature} & \textbf{Protobuf} & \textbf{JSON} \\
\hline
Message Size & 17 bytes (example) & 65 bytes \\
Serialization & 1.2 $\mu$s & 4.5 $\mu$s \\
Type Safety & Compile-time & Runtime \\
Schema Evolution & Built-in & Manual \\
Code Generation & Automatic & Manual \\
\hline
\end{tabular}
\caption{Protobuf vs JSON Comparison}
\end{table}

\subsection{Architecture}

\begin{figure}[H]
\centering
\begin{tikzpicture}[
    node distance=1.5cm,
    box/.style={draw, rectangle, rounded corners, minimum width=2.5cm, minimum height=1cm, align=center, font=\small},
    arrow/.style={->, thick}
]
    % Proto source
    \node[box, fill=yellow!30] (proto) {star-proto\\(.proto files)};

    % Generators
    \node[box, fill=blue!20, below left=2cm and 1cm of proto] (kotlin) {protoc\\(Kotlin)};
    \node[box, fill=green!20, below=2cm of proto] (ts) {protoc\\(TypeScript)};
    \node[box, fill=orange!20, below right=2cm and 1cm of proto] (c) {nanopb\\(C)};

    % Output
    \node[box, fill=blue!10, below=1.5cm of kotlin] (gateway) {star-gateway\\(*.kt)};
    \node[box, fill=green!10, below=1.5cm of ts] (ui) {star-ui\\(*\_pb.ts)};
    \node[box, fill=orange!10, below=1.5cm of c] (firmware) {star-esp32-firmware\\(*.pb.c)};

    % Connections
    \draw[arrow] (proto) -- (kotlin);
    \draw[arrow] (proto) -- (ts);
    \draw[arrow] (proto) -- (c);
    \draw[arrow] (kotlin) -- (gateway);
    \draw[arrow] (ts) -- (ui);
    \draw[arrow] (c) -- (firmware);
\end{tikzpicture}
\caption{Code Generation Flow}
\end{figure}

% =============================================================================
\section{Directory Structure}
% =============================================================================

\begin{lstlisting}[language=bash]
star-proto/
|-- STYLE_GUIDE.md                # Boston Dynamics-based style guide
|-- buf.gen.yaml                  # Code generation config
|
|-- proto/
|   |-- buf.yaml                  # Buf lint/build configuration
|   +-- star/
|       +-- v1/                   # Versioned package (star.v1)
|           |-- common.proto      # Shared types (headers, status, vectors)
|           |-- motor_control.proto  # Motor command service
|           +-- telemetry.proto   # Telemetry streaming service
|
|-- nanopb/                       # nanopb options for ESP32 constraints
|   |-- common.options
|   |-- motor_control.options
|   +-- telemetry.options
|
|-- tests/                        # Serialization tests
|   |-- go/                       # Go tests (primary)
|   |-- typescript/               # TypeScript tests
|   +-- nanopb/                   # C/nanopb tests
|
+-- gen/                          # Generated code (gitignored)
    |-- go/                       # Go code for star-gateway
    |-- typescript/               # TypeScript for star-ui
    +-- nanopb/                   # C code for ESP32
\end{lstlisting}

% =============================================================================
\section{Protocol Definitions}
% =============================================================================

\subsection{motor\_control.proto}

\begin{lstlisting}[language=Protobuf, caption=star/v1/motor\_control.proto]
syntax = "proto3";

package star.v1;

import "star/v1/common.proto";

option go_package = "github.com/Locked-Inc/star-proto/gen/go/star/v1;starv1";
option java_multiple_files = true;
option java_outer_classname = "MotorControlProto";
option java_package = "com.star.proto.v1";

// Motor control service for low-latency differential drive commands.
service MotorControlService {
    // Set wheel velocities for differential drive.
    rpc SetVelocity(SetVelocityRequest) returns (SetVelocityResponse);

    // Immediate emergency stop. Stops all motors and engages brakes.
    rpc EmergencyStop(EmergencyStopRequest) returns (EmergencyStopResponse);

    // Set individual motor power directly (bypass PID).
    rpc SetMotorPower(SetMotorPowerRequest) returns (SetMotorPowerResponse);

    // Stream encoder data from all motors.
    rpc StreamEncoders(StreamEncodersRequest) returns (stream EncoderData);

    // Bidirectional control stream for continuous control.
    rpc ControlStream(stream VelocityCommand) returns (stream EncoderData);
}

// Request to set wheel velocities.
message SetVelocityRequest {
    // Standard request header.
    RequestHeader header = 1;

    // Velocity command to execute.
    VelocityCommand command = 2;
}

// Response to velocity command.
message SetVelocityResponse {
    // Standard response header with status.
    ResponseHeader header = 1;

    // Current motor status after command.
    repeated MotorStatus motor_status = 2;
}

// Velocity command for differential drive robot.
// Uses MKS units: meters per second.
message VelocityCommand {
    // Left wheel velocity in meters per second.
    // Valid range: -2.0 to 2.0 m/s.
    double left_velocity_mps = 1;

    // Right wheel velocity in meters per second.
    // Valid range: -2.0 to 2.0 m/s.
    double right_velocity_mps = 2;

    // Command sequence number for ordering.
    uint32 sequence = 3;

    // Client timestamp in microseconds.
    int64 timestamp_us = 4;
}

// Real-time encoder feedback from motor controller.
message EncoderData {
    // Motor index: 0 = left, 1 = right.
    int32 motor_id = 1;

    // Absolute encoder position in ticks.
    int64 ticks = 2;

    // Current velocity in meters per second.
    double velocity_mps = 3;

    // Timestamp in microseconds since boot.
    int64 timestamp_us = 4;
}

// Motor operating state.
enum MotorState {
    // Unknown state (zero value per Boston Dynamics style).
    MOTOR_STATE_UNKNOWN = 0;
    MOTOR_STATE_IDLE = 1;
    MOTOR_STATE_RUNNING = 2;
    MOTOR_STATE_FAULT = 3;
    MOTOR_STATE_ESTOP = 4;
}
\end{lstlisting}

\subsection{telemetry.proto}

\begin{lstlisting}[language=Protobuf, caption=star/v1/telemetry.proto]
syntax = "proto3";

package star.v1;

import "google/protobuf/timestamp.proto";
import "star/v1/common.proto";

option go_package = "github.com/Locked-Inc/star-proto/gen/go/star/v1;starv1";
option java_multiple_files = true;
option java_outer_classname = "TelemetryProto";
option java_package = "com.star.proto.v1";

// Telemetry streaming service for real-time sensor data.
service TelemetryService {
    // Get current telemetry snapshot.
    rpc GetTelemetry(GetTelemetryRequest) returns (GetTelemetryResponse);

    // Stream telemetry at specified rate.
    rpc StreamTelemetry(StreamTelemetryRequest) returns (stream TelemetryData);

    // Get system status information.
    rpc GetSystemStatus(GetSystemStatusRequest) returns (GetSystemStatusResponse);
}

// Request for telemetry snapshot.
message GetTelemetryRequest {
    RequestHeader header = 1;
}

// Response with telemetry snapshot.
message GetTelemetryResponse {
    ResponseHeader header = 1;
    TelemetryData telemetry = 2;
}

// Complete telemetry snapshot from all sensors.
message TelemetryData {
    ImuData imu = 1;
    GpsData gps = 2;
    double battery_percent = 3;       // Range: 0 to 100
    int32 wifi_signal_dbm = 4;        // Typical: -100 to -30 dBm
    double cpu_usage_percent = 5;     // Range: 0 to 100
    double temperature_celsius = 6;   // From DS18B20 sensor
    double motor_load_percent = 7;    // Average across motors
    google.protobuf.Timestamp timestamp = 8;
}

// IMU sensor data with MKS units.
message ImuData {
    double pitch_rad = 1;             // Range: -pi/2 to pi/2
    double roll_rad = 2;              // Range: -pi to pi
    double yaw_rad = 3;               // Range: -pi to pi
    double accel_x_mps2 = 4;          // Linear acceleration X (m/s^2)
    double accel_y_mps2 = 5;          // Linear acceleration Y (m/s^2)
    double accel_z_mps2 = 6;          // ~9.81 at rest
    double gyro_x_rad_per_s = 7;      // Angular velocity X (rad/s)
    double gyro_y_rad_per_s = 8;      // Angular velocity Y (rad/s)
    double gyro_z_rad_per_s = 9;      // Angular velocity Z (rad/s)
}

// GPS position data.
message GpsData {
    double latitude_deg = 1;          // WGS84, range: -90 to 90
    double longitude_deg = 2;         // WGS84, range: -180 to 180
    double altitude_m = 3;            // Meters above sea level
    double accuracy_m = 4;            // Horizontal accuracy (m)
    int32 satellites = 5;
    GpsFix fix_type = 6;
}

// GPS fix type enumeration.
enum GpsFix {
    GPS_FIX_UNKNOWN = 0;              // Zero value is UNKNOWN
    GPS_FIX_NONE = 1;
    GPS_FIX_2D = 2;
    GPS_FIX_3D = 3;
    GPS_FIX_DGPS = 4;
    GPS_FIX_RTK_FLOAT = 5;
    GPS_FIX_RTK_FIXED = 6;
}

// System status information.
message SystemStatus {
    ConnectionStatus connection_status = 1;
    RobotMode mode = 2;
    bool esp32_connected = 3;
    bool lidar_connected = 4;
    bool ros_connected = 5;
    string firmware_version = 6;
    uint64 uptime_s = 7;
    uint32 free_heap_bytes = 8;
}

// Connection state enumeration.
enum ConnectionStatus {
    CONNECTION_STATUS_UNKNOWN = 0;    // Zero value is UNKNOWN
    CONNECTION_STATUS_CONNECTED = 1;
    CONNECTION_STATUS_DISCONNECTED = 2;
    CONNECTION_STATUS_CONNECTING = 3;
    CONNECTION_STATUS_ERROR = 4;
}

// Robot operating mode enumeration.
enum RobotMode {
    ROBOT_MODE_UNKNOWN = 0;           // Zero value is UNKNOWN
    ROBOT_MODE_IDLE = 1;
    ROBOT_MODE_MANUAL = 2;
    ROBOT_MODE_AUTONOMOUS = 3;
    ROBOT_MODE_MAPPING = 4;
    ROBOT_MODE_EMERGENCY_STOP = 5;
}
\end{lstlisting}

\subsection{navigation.proto (Planned)}

\begin{lstlisting}[language=Protobuf, caption=star/v1/navigation.proto (planned)]
syntax = "proto3";

package star.v1;

import "google/protobuf/timestamp.proto";
import "star/v1/common.proto";

option go_package = "github.com/Locked-Inc/star-proto/gen/go/star/v1;starv1";
option java_multiple_files = true;
option java_outer_classname = "NavigationProto";
option java_package = "com.star.proto.v1";

// Navigation service for autonomous path planning and execution.
service NavigationService {
    // Get current navigation state.
    rpc GetState(GetNavigationStateRequest) returns (GetNavigationStateResponse);

    // Add waypoint to navigation queue.
    rpc AddWaypoint(AddWaypointRequest) returns (AddWaypointResponse);

    // Clear all queued waypoints.
    rpc ClearWaypoints(ClearWaypointsRequest) returns (ClearWaypointsResponse);

    // Start autonomous navigation to queued waypoints.
    rpc StartNavigation(StartNavigationRequest) returns (StartNavigationResponse);

    // Stop navigation and hold position.
    rpc StopNavigation(StopNavigationRequest) returns (StopNavigationResponse);

    // Stream navigation updates.
    rpc StreamNavigation(StreamNavigationRequest) returns (stream NavigationState);

    // Stream LiDAR scan data.
    rpc StreamLidar(StreamLidarRequest) returns (stream LidarScan);
}

// Current navigation state.
message NavigationState {
    SE2Pose current_pose = 1;         // Current robot pose
    repeated SE2Pose path = 2;        // Planned path points
    repeated Waypoint waypoints = 3;  // Queued waypoints
    bool is_navigating = 4;
    double progress_percent = 5;      // 0-100
    double eta_s = 6;                 // Estimated time remaining
    NavigationStatus status = 7;
}

// Navigation waypoint.
message Waypoint {
    string id = 1;
    double x_m = 2;                   // X position in meters
    double y_m = 3;                   // Y position in meters
    string label = 4;
    WaypointType type = 5;
}

// Waypoint type enumeration.
enum WaypointType {
    WAYPOINT_TYPE_UNKNOWN = 0;        // Zero value is UNKNOWN
    WAYPOINT_TYPE_NORMAL = 1;
    WAYPOINT_TYPE_STOP = 2;
    WAYPOINT_TYPE_SLOW = 3;
}

// Navigation status enumeration.
enum NavigationStatus {
    NAVIGATION_STATUS_UNKNOWN = 0;    // Zero value is UNKNOWN
    NAVIGATION_STATUS_IDLE = 1;
    NAVIGATION_STATUS_NAVIGATING = 2;
    NAVIGATION_STATUS_PAUSED = 3;
    NAVIGATION_STATUS_GOAL_REACHED = 4;
    NAVIGATION_STATUS_ERROR = 5;
}

// LiDAR scan data.
message LidarScan {
    repeated LidarPoint points = 1;
    google.protobuf.Timestamp timestamp = 2;
    string frame_id = 3;
}

// Single LiDAR measurement point.
message LidarPoint {
    double angle_rad = 1;             // Angle in radians
    double distance_m = 2;            // Distance in meters
    double intensity = 3;             // Intensity 0-255
}

// Occupancy grid map.
message OccupancyGrid {
    double resolution_m = 1;          // Meters per cell
    uint32 width = 2;
    uint32 height = 3;
    SE2Pose origin = 4;
    bytes data = 5;                   // -1=unknown, 0=free, 100=occupied
    google.protobuf.Timestamp timestamp = 6;
}
\end{lstlisting}

\subsection{common.proto}

\begin{lstlisting}[language=Protobuf, caption=star/v1/common.proto]
syntax = "proto3";

package star.v1;

import "google/protobuf/duration.proto";
import "google/protobuf/timestamp.proto";

option go_package = "github.com/Locked-Inc/star-proto/gen/go/star/v1;starv1";
option java_multiple_files = true;
option java_package = "com.star.proto.v1";

// Standard request header included in all requests.
message RequestHeader {
    string request_id = 1;                    // UUID for tracing
    google.protobuf.Timestamp client_timestamp = 2;
    string client_version = 3;                // e.g., "1.0.0"
    google.protobuf.Duration timeout = 4;     // Optional deadline
}

// Standard response header included in all responses.
message ResponseHeader {
    string request_id = 1;                    // Echo of request_id
    google.protobuf.Timestamp server_timestamp = 2;
    Status status = 3;                        // Processing status
    string error_message = 4;                 // If status != STATUS_OK
    int64 latency_us = 5;                     // Round-trip latency
}

// Response status codes (zero value is UNKNOWN per Boston Dynamics).
enum Status {
    STATUS_UNKNOWN = 0;
    STATUS_OK = 1;
    STATUS_INVALID_REQUEST = 2;
    STATUS_INTERNAL_ERROR = 3;
    STATUS_NOT_FOUND = 4;
    STATUS_TIMEOUT = 5;
    STATUS_INVALID_STATE = 6;
    STATUS_ESTOP_ACTIVE = 7;
}

// 3D vector with MKS units (context determines: m, m/s, etc).
message Vector3 {
    double x = 1;
    double y = 2;
    double z = 3;
}

// Quaternion for orientation (w, x, y, z).
message Quaternion {
    double w = 1;
    double x = 2;
    double y = 3;
    double z = 4;
}

// SE2 pose for 2D navigation.
message SE2Pose {
    double x_m = 1;           // X position in meters
    double y_m = 2;           // Y position in meters
    double angle_rad = 3;     // Heading in radians (-pi to pi)
}

// SE2 velocity for 2D motion.
message SE2Velocity {
    double vx_mps = 1;        // Linear velocity X (m/s)
    double vy_mps = 2;        // Linear velocity Y (m/s)
    double omega_rad_per_s = 3; // Angular velocity (rad/s)
}
\end{lstlisting}

% =============================================================================
\section{Code Generation}
% =============================================================================

\subsection{Prerequisites}

\begin{lstlisting}[language=bash]
# Install protoc compiler
# macOS
brew install protobuf

# Ubuntu
sudo apt install protobuf-compiler

# Install Buf (recommended)
# macOS
brew install bufbuild/buf/buf

# Other platforms
npm install -g @bufbuild/buf
\end{lstlisting}

\subsection{Using Buf (Recommended)}

Create \texttt{proto/buf.yaml}:

\begin{lstlisting}[language=yaml, caption=proto/buf.yaml]
version: v1
name: buf.build/star/proto

deps:
  - buf.build/googleapis/googleapis

lint:
  use:
    - DEFAULT
    - COMMENTS
  except:
    - PACKAGE_VERSION_SUFFIX
    # Boston Dynamics uses _UNKNOWN, not _UNSPECIFIED
    - ENUM_ZERO_VALUE_SUFFIX
    - RPC_REQUEST_RESPONSE_UNIQUE
    - RPC_REQUEST_STANDARD_NAME
    - RPC_RESPONSE_STANDARD_NAME
  service_suffix: Service
  rpc_allow_same_request_response: true

breaking:
  use:
    - FILE
    - PACKAGE
\end{lstlisting}

Create \texttt{buf.gen.yaml}:

\begin{lstlisting}[language=yaml, caption=buf.gen.yaml]
version: v1

managed:
  enabled: true
  go_package_prefix:
    default: github.com/Locked-Inc/star-proto/gen/go
  optimize_for: SPEED

plugins:
  # Go for star-gateway (primary language)
  - plugin: buf.build/protocolbuffers/go
    out: gen/go
    opt:
      - paths=source_relative

  - plugin: buf.build/grpc/go
    out: gen/go
    opt:
      - paths=source_relative

  # TypeScript for star-ui (protobuf-ts)
  - plugin: buf.build/community/timostamm-protobuf-ts
    out: gen/typescript
    opt:
      - long_type_string
      - optimize_code_size
      - generate_dependencies
      - ts_nocheck
      - eslint_disable

  # C for star-esp32-firmware (nanopb)
  # Note: nanopb requires local installation: pip install nanopb
  # - plugin: nanopb_generator
  #   out: gen/nanopb
  #   opt:
  #     - --options-path=nanopb/
\end{lstlisting}

Generate code:

\begin{lstlisting}[language=bash]
# From star-proto directory
cd star-proto

# Lint proto files
buf lint proto/

# Format proto files
buf format --diff proto/

# Generate all targets
buf generate proto/ --template buf.gen.yaml --include-imports
\end{lstlisting}

\subsection{Manual Generation Scripts (Alternative to Buf)}

\texttt{scripts/generate-go.sh} (Primary):

\begin{lstlisting}[language=bash, caption=generate-go.sh]
#!/bin/bash
set -e

PROTO_DIR="proto"
OUT_DIR="gen/go"

mkdir -p $OUT_DIR

protoc \
    -I=$PROTO_DIR \
    --go_out=$OUT_DIR \
    --go_opt=paths=source_relative \
    --go-grpc_out=$OUT_DIR \
    --go-grpc_opt=paths=source_relative \
    $PROTO_DIR/star/v1/*.proto

echo "Go code generated in $OUT_DIR"
\end{lstlisting}

\texttt{scripts/generate-nanopb.sh} (ESP32):

\begin{lstlisting}[language=bash, caption=generate-nanopb.sh]
#!/bin/bash
set -e

PROTO_DIR="proto"
OUT_DIR="gen/nanopb"

mkdir -p $OUT_DIR

for proto in $PROTO_DIR/star/v1/*.proto; do
    nanopb_generator \
        --proto-path=$PROTO_DIR \
        --options-path=nanopb \
        --output-dir=$OUT_DIR \
        "$proto"
done

echo "nanopb C code generated in $OUT_DIR"
\end{lstlisting}

% =============================================================================
\section{Best Practices}
% =============================================================================

\subsection{Schema Evolution}

\begin{enumerate}
    \item \textbf{Never remove fields} - mark as \texttt{reserved}
    \item \textbf{Never change field numbers} - they define wire format
    \item \textbf{Add new fields} with new numbers only
    \item \textbf{Use optional} for fields that may not be present
\end{enumerate}

\begin{lstlisting}[language=Protobuf, caption=Schema Evolution Example]
message VelocityCommand {
    double left_velocity_mps = 1;
    double right_velocity_mps = 2;

    // Reserved: removed fields
    reserved 3;
    reserved "old_field_name";

    // New optional field (added later)
    optional double acceleration_limit_mps2 = 4;
}
\end{lstlisting}

\subsection{Naming Conventions (Boston Dynamics Style)}

\begin{table}[H]
\centering
\begin{tabular}{|l|l|l|}
\hline
\textbf{Element} & \textbf{Convention} & \textbf{Example} \\
\hline
Package & lowercase.version & star.v1 \\
Message & PascalCase & VelocityCommand \\
Field & snake\_case\_with\_units & left\_velocity\_mps \\
Enum type & PascalCase & RobotMode \\
Enum value & TYPE\_PREFIX\_VALUE & ROBOT\_MODE\_MANUAL \\
Enum zero & TYPE\_PREFIX\_UNKNOWN & ROBOT\_MODE\_UNKNOWN \\
Service & PascalCaseService & MotorControlService \\
RPC method & PascalCase & SetVelocity \\
\hline
\end{tabular}
\caption{Protobuf Naming Conventions (Boston Dynamics Style)}
\end{table}

\textbf{Key Boston Dynamics Conventions:}
\begin{itemize}
    \item \textbf{Enum zero value}: Always ends with \texttt{\_UNKNOWN}, not \texttt{\_UNSPECIFIED}
    \item \textbf{Enum prefixing}: Values prefixed with type name (e.g., \texttt{MOTOR\_STATE\_IDLE})
    \item \textbf{Unit suffixes}: Fields include units (e.g., \texttt{\_mps}, \texttt{\_rad}, \texttt{\_celsius}, \texttt{\_us})
    \item \textbf{Double precision}: Use \texttt{double} instead of \texttt{float} for numeric values
    \item \textbf{Request/Response headers}: All RPCs include \texttt{RequestHeader}/\texttt{ResponseHeader}
\end{itemize}

\subsection{Documentation}

Always document messages and fields with units in field names:

\begin{lstlisting}[language=Protobuf]
// Velocity command for differential drive robot.
// Sent from controller to gateway at up to 100Hz.
message VelocityCommand {
    // Left wheel velocity in meters per second.
    // Valid range: -2.0 to 2.0 m/s.
    // Positive = forward, negative = reverse.
    double left_velocity_mps = 1;

    // Right wheel velocity in meters per second.
    // Valid range: -2.0 to 2.0 m/s.
    // Positive = forward, negative = reverse.
    double right_velocity_mps = 2;

    // Command sequence number for ordering.
    uint32 sequence = 3;

    // Client timestamp in microseconds.
    int64 timestamp_us = 4;
}
\end{lstlisting}

% =============================================================================
\section{Integration}
% =============================================================================

\subsection{star-gateway Integration (Go)}

Add to \texttt{star-gateway/go.mod}:

\begin{lstlisting}[language=bash]
module github.com/Locked-Inc/star-gateway

go 1.21

require (
    github.com/Locked-Inc/star-proto/gen/go v0.0.0
    google.golang.org/grpc v1.60.0
    google.golang.org/protobuf v1.31.0
)

// For local development, use replace directive:
replace github.com/Locked-Inc/star-proto/gen/go => ../star-proto/gen/go
\end{lstlisting}

Usage in Go code:

\begin{lstlisting}[language=Go]
import (
    starv1 "github.com/Locked-Inc/star-proto/gen/go/star/v1"
    "google.golang.org/grpc"
)

// Create velocity command
cmd := &starv1.VelocityCommand{
    LeftVelocityMps:  1.5,
    RightVelocityMps: 1.5,
    Sequence:         42,
}
\end{lstlisting}

\subsection{star-ui Integration (TypeScript)}

Add npm script to \texttt{star-ui/package.json}:

\begin{lstlisting}[language=json]
{
  "scripts": {
    "proto:generate": "cd ../star-proto && buf generate proto/"
  },
  "dependencies": {
    "@protobuf-ts/runtime": "^2.9.0",
    "@protobuf-ts/grpcweb-transport": "^2.9.0"
  }
}
\end{lstlisting}

\subsection{star-esp32-firmware Integration (nanopb)}

Copy generated files to firmware:

\begin{lstlisting}[language=bash]
# After running buf generate or nanopb_generator
cp star-proto/gen/nanopb/*.pb.c star-esp32-firmware/lib/star_proto/src/
cp star-proto/gen/nanopb/*.pb.h star-esp32-firmware/lib/star_proto/include/

# Include in platformio.ini
lib_deps =
    nanopb/Nanopb@^0.4.8
\end{lstlisting}

% =============================================================================
\section{References}
% =============================================================================

\begin{itemize}
    \item Boston Dynamics Protobuf Style Guide: \url{https://dev.bostondynamics.com/docs/protos/style_guide}
    \item Protocol Buffers Documentation: \url{https://protobuf.dev/}
    \item gRPC Documentation: \url{https://grpc.io/docs/}
    \item Buf CLI: \url{https://buf.build/docs/}
    \item gRPC-Web: \url{https://github.com/grpc/grpc-web}
    \item nanopb (C): \url{https://jpa.kapsi.fi/nanopb/}
    \item protobuf-ts (TypeScript): \url{https://github.com/timostamm/protobuf-ts}
    \item STAR Software Architecture Primer: Section 3 (Protocol Analysis)
\end{itemize}

